% !TeX root = ../packet-and-principal-models.tex
\begin{thebibliography}{99}
\bibitem{DeligneHodge2006}
P. Deligne,
\emph{The Hodge Conjecture}, in J. Carlson, A. Jaffe and A. Wiles (eds.),
\emph{The Millennium Prize Problems}, Clay Mathematics Institute/AMS,
2006, 45--53.

\bibitem{Goodson2024}
H. Goodson,
\emph{An exploration of degeneracy in abelian varieties of Fermat type},
arXiv:2211.03909v2, 1 August 2024.
The cited numbering refers to this version.
\url{https://arxiv.org/abs/2211.03909v2}

\bibitem{LiMantovanPries2021}
W.~Li, E.~Mantovan and R.~Pries,
\emph{Data for Shimura varieties intersecting the Torelli locus},
arXiv:2105.02286v1 (2021); Pacific J. Math. \textbf{343} (2026), 179--210.
Proposition~4.5, Proposition~6.11 and Lemma~6.12 refer to the
inspected arXiv version.
\url{https://arxiv.org/abs/2105.02286v1}

\bibitem{Markman2025}
E. Markman,
\emph{Cycles on abelian $2n$-folds of Weil type from secant sheaves on abelian $n$-folds},
arXiv:2502.03415v2 (2025), Section~1.1, p.~3, and Theorem~1.5.1, p.~9.
The theorem proves algebraicity of Hodge--Weil classes for polarized
abelian sixfolds of Weil type with discriminant $-1$, without a
principal-polarization hypothesis.
\url{https://arxiv.org/pdf/2502.03415v2}

\bibitem{Masley1976}
J. M. Masley,
Solution of small class number problems for cyclotomic fields,
\emph{Compositio Mathematica} \textbf{33} (1976), 179--186,
\url{https://www.numdam.org/item/CM_1976__33_2_179_0.pdf}.

\bibitem{McMullen2026}
C. T. McMullen,
\emph{Triangle groups and Hilbert modular varieties},
J. Math. Soc. Japan 78 (2026), no.~3, 729--749; DOI 10.2969/jmsj/94179417.
The theorem and section locators used here refer to the author's version
dated 27 June 2023, typeset 27 March 2026, Sections~4 and~6
(printed pp.~14--15 and 19--21).
\url{https://people.math.harvard.edu/~ctm/papers/home/text/papers/quat/quat.pdf}

\bibitem{MilneModuli2011}
J. S. Milne,
\emph{Shimura Varieties and Moduli}, manuscript dated 30 April 2011,
Theorems 3.6 and 7.13, pp.~15 and 38--39. The numbering here refers to
this author version.
\url{https://www.jmilne.org/math/xnotes/svh.pdf}

\bibitem{MilneCM2020}
J. S. Milne,
\emph{Complex Multiplication}, version 0.10, 14 July 2020,
Section 3.11 and Propositions 3.12--3.13 and 3.17, pp.~29--31.
\url{https://www.jmilne.org/math/CourseNotes/CM.pdf}

\bibitem{MilneHodge2020}
J. S. Milne,
\emph{Hodge classes on abelian varieties},
arXiv:2010.08857, 2020.
\url{https://arxiv.org/abs/2010.08857}

\bibitem{Mostaed2026}
A. Mostaed,
\emph{McMullen's Curve, the Weil Locus, and the Hodge Conjecture for Abelian Sixfolds},
arXiv:2603.20268v1, 14 March 2026.
\url{https://arxiv.org/abs/2603.20268}

\bibitem{OpenAICM2026}
OpenAI math repository,
\emph{The rational Hodge conjecture for CM abelian varieties},
public manuscript dated 6 October 2026, introduction (main CM theorem)
and tensor-reduction section (auxiliary-factor contractions).
The comparison uses repository commit
\texttt{4eaee87548250b2682affda7d05d947c05554bae};
the general theorem and its formalization are not proof inputs here.
\url{https://github.com/openai/math/tree/4eaee87548250b2682affda7d05d947c05554bae/preprints/The-rational-Hodge-conjecture-for-CM-abelian-varieties-October-6-2026}

\bibitem{Saller2026}
K. Saller,
\emph{Split Weil structures on abelian sixfolds and the Hodge conjecture},
version 2, Zenodo record 23290280 (2026), Remarks~4.7--4.8.
\url{https://doi.org/10.5281/zenodo.23290280}

\bibitem{Schoen1998}
C. Schoen,
\emph{Addendum to ``Hodge classes on self-products of a variety with an automorphism''},
Compositio Math. 114 (1998), no.~3, 329--336.

\end{thebibliography}
